% !TeX root = ../main.tex
\subsection{现有实现的补充验证}\label{minor:supplement}
在原有规划结构上，后续实现进一步处理建筑热状态、设备工况、负荷压缩与网络物理反馈。表\ref{tab:minor-evidence}汇总已经保存的独立实验，均保持各自的输入对象与边界，不将不同实验的改善率合并为同一工程收益。

\begin{table}[H]
\centering
\caption{保留原规划框架后补入的验证记录}\label{tab:minor-evidence}
\begin{tabularx}{\linewidth}{@{}>{\raggedright\arraybackslash}L{0.22\linewidth}YY@{}}
\toprule
模块与对象 & 验证条件 & 已有结果\\\midrule
代表日权重\newline 六栋BDG2计量建筑 & 固定8--10个实际代表日，以日数与能量等式约束权重 & 归一$W_1$误差中位数由12.3150\%降至4.4775\%；完整统计准入由3/6增至4/6\\
厂家工况调度\newline 三城、两型号 & 同型号165对制热量／输入功率；同资产、同服务，2880小时 & 电量费降低1.10\%--1.60\%；同时记录增加的用电量\\
连续运行\newline 武汉四楼数学场景 & 固定规划投资；8784小时服务年及2928小时声明冬季 & 两组建筑状态检查均为零违约；运行gap分别0.1344\%、0.4755\%\\
物理反馈\newline 72小时参考管网 & 同一交付量；按开启工况检查温度、压降并反馈排除模式 & 第三轮最大压降237.96~Pa/m、远端供水57.45~$^\circ$C，满足300~Pa/m与55~$^\circ$C阈值\\
\bottomrule
\end{tabularx}
\end{table}

\textbf{守恒与设备边界。}负荷基准采用BDG2公开计量\supcite{29}；权重校正保持总日数与总能量，直接减小分布差异。厂家实验按TICA同型号有效工况\supcite{30}分别插值制热量$Q$和输入功率$P$后求$\mathrm{COP}$，空白角点保持缺失。三城采用同一2880小时窗口的公开天气\supcite{31}与共同计费条件，隔离气候和调度影响。该组节费来自时移，增加用电另行记录。

\textbf{连续状态与工程反馈。}四楼回放采用已登记合成参数，服务年为北京时间2023年7月1日至2024年7月1日，声明冬季为2023年11月15日至2024年3月16日，均含起点、不含终点。建筑采用精确1R1C与固定初末室温，独立复算最大状态残差分别为$1.10\times10^{-13}$与$6.75\times10^{-14}$~$^\circ$C。另在独立泵星型支路的有限稳态模式中，每轮将不合格模式禁用后重新优化；三轮交付热量均为191.5475~kWh，达标可行方案的参考费用由178.64增至224.26元，记录了物理约束带来的成本。

上述记录分别对应随附证据索引E01、E02、E08、E04。本次文本修订沿用已保存的运行结果；原始聚合算例与新增逐楼实验分别定位，便于按模型版本复算。
